% Independent audits: root and descent_groups; certificate code/C_finish_torsion_orders.py.
\begin{theorem}[Proved: torsion of full ramification height gives infinite order]
\label{C:ramified_torsion_infinite}
Let $A$ be any characteristic-zero DVR of residue characteristic
$p$, let $\pi$ be a uniformizer, and write $p=u\pi^e$ with $u$
a unit. Let $J$ be any $A$-module containing an element $x$ of
exact $\pi$-order $e$:
\[
 \pi^e x=0,\qquad \pi^{e-1}x\ne0.
\]
For the free line $L_x$ over $k=A\ltimes J$ on which the chosen
generator of $C_p$ acts by $1+x$, the class
\[
 \xi_x=[\operatorname{Ind}_{C_p}^{\mathfrak S_p}L_x]
        -[\operatorname{Ind}_{C_p}^{\mathfrak S_p}\mathbf1]
\]
has infinite additive order in the cokernel of the degree-$p$
polynomial map $\vartheta_k$. It vanishes after inverting $p$
and at every residue-field point. This holds also for $p=2$.
\end{theorem}
\begin{proof}
Apply the primary curvature-length invariant of
Lemma~\ref{C:curvature_length} to $R=A/\pi^m$, with
$C^i=\widetilde d^{i+1}\widetilde d^i/\pi^m$ and
$q=z-\bar u$.

Fix any finitely generated $A$-submodule $K\subseteq J$ containing
$x$. The DVR structure theorem gives a cyclic summand $A/\pi^r$,
with $r\ge e$, on which the projection of $x$ has exact order $e$.
After multiplication by a unit, projection therefore gives
$\phi:K\to A/\pi^r$ with $\phi(x)=\pi^{r-e}$.
For every integer $m\ge2r$ there is a ring map
\[
 A\ltimes K\longrightarrow A/\pi^m,\qquad
 (a,j)\longmapsto a+\pi^{m-r}\phi(j).
                                                        \tag{C.1.30}
\]
It is well-defined because $\pi^r\phi(j)=0$, and multiplicative
because $2(m-r)\ge m$. It sends $x$ to $\pi^{m-e}$.

Choose additionally $m>2e$. The scalar
$T=1+\pi^{m-e}$ has $T^p=1$ modulo $\pi^m$.
Its lifted Tate curvature satisfies
\[
 \frac{T^p-1}{\pi^m}\equiv u\pmod\pi.
\]
Indeed, for $2\le i<p$ the $i$-th binomial term after division
has valuation at least $(i-1)(m-e)>0$; the last term has
valuation $(p-1)m-pe>0$. For $p=2$ only that last term occurs,
and its valuation is $m-2e>0$.
The full scalar Tate complex is therefore on its
$(z-\bar u)$-primary part, and
\[
 \nu_{z-\bar u,m}(L_x)=m-e=-e\pmod m.
\]
For each generator twist $g\mapsto g^a$, $1\le a<p$, the same
calculation gives residual curvature $a\bar u$ and degree-zero
cokernel length $m-e$. Exactly $a=1$ contributes to this primary
invariant. Mackey restriction of the induced module gives each
normalizer twist once; all other terms are regular actions that
lift to $A$. Hence, after the scalar extension (C.1.30), the
restricted induced difference has invariant $-e$ modulo $m$.
Every Young permutation class lifts to $A$ and has invariant zero.

If a positive multiple $N\xi_x$ belonged to the polynomial image
over $A\ltimes J$, filtered-colimit continuity
(Proposition~\ref{C:filtered_colimits}) would realize that equality
over some finite $K$ containing $x$. Here the local rings
$A\ltimes K$ have $K_0=\mathbb Z$, so the polynomial image is
generated by the same constant Young permutation classes.
Choose $m>\max(2r,2e,Ne)$ in (C.1.30). The invariant of this
alleged equality is $-Ne\not\equiv0\pmod m$, a contradiction.
Finally $x$ is $p$-torsion and square-zero, giving the asserted
vanishing at the generic and residue-field points.
\end{proof}

\begin{corollary}[Proved: every unramified torsion witness has infinite order]
\label{C:unramified_torsion_infinite}
If $A$ is a characteristic-zero DVR with maximal ideal $(p)$,
for any prime $p$, and $0\ne x\in J[p]$, then the class $\xi_x$
of Theorem~\ref{C:ramified_torsion_infinite} has infinite order in
the degree-$p$ polynomial cokernel. Thus the odd-prime
nonvanishing of Proposition~\ref{C:odd_torsion_allheights} can be
strengthened to infinite order, and the same holds for the
rank-one dyadic torsion witness. Over an arbitrary ramified DVR,
any module with nonzero $p$-divisible torsion contains a witness
of Theorem~\ref{C:ramified_torsion_infinite}.
\end{corollary}
\begin{proof}
In the unramified case $e=1$, so every nonzero element of $J[p]$
has exactly the required order. For the last assertion choose a
nonzero torsion element of a $p$-divisible torsion submodule.
Successive preimages under $p$ have arbitrarily large finite
$\pi$-orders; multiplication by a suitable power of $\pi$ then
produces exact order $e$.
\end{proof}
